\documentclass[conference]{IEEEtran}
\usepackage{booktabs}
\usepackage{amsmath}
\usepackage{url}
\usepackage{array}
\usepackage[hidelinks]{hyperref}

\begin{document}

\title{Does a Small LLM Read Purchase Timing?\\Zero-Shot Reranking of Customers' Own Past Assets on FAR-Trans}

\author{\IEEEauthorblockN{Mark Starchenko}
\IEEEauthorblockA{Team ``name'' (A.~Kosychev, K.~Patrushev, M.~Starchenko, N.~Vaal)\\
HSE University, Graduate School of Business, mastarchenko@edu.hse.ru\\
Session 2 report, LLM4Rec, 25 November 2026\\
Repository: \url{https://github.com/ak1232320/rec-sys-hw1}, branch \texttt{starchenko/session2-prefix-rerank},\\
folder \texttt{hw-module-2-2/starchenko}, experiment commit \texttt{bb64f11}}}

\maketitle

\begin{abstract}
I implemented the Ranker question from Session~1: Qwen3-1.7B, prompted zero-shot with dated Buy and Sell events, reorders only the assets a FAR-Trans customer has already bought, and the popularity fill of \texttt{latest\_buy} stays fixed. On the starter's 100-user sample (cutoff 2021-06-01, 30 days, 11 users with purchases) it raises nDCG@10 from 0.0277 for \texttt{latest\_buy} to 0.0316, below a logistic-regression ranker that receives the same facts (0.0343). A control without any model, which keeps the order in which the prompt lists the assets, scores 0.0341, and 46 of the 57 valid LLM answers return exactly that order, so the gain does not come from reading dates. Three of 60 answers were invalid and fell back to recency; the sample is too small to rank the methods, and the next experiment tests order bias on the nine-cutoff development protocol.
\end{abstract}

\begin{IEEEkeywords}
financial asset recommendation, repeat consumption, LLM reranking, position bias, FAR-Trans
\end{IEEEkeywords}

\section{Introduction}

\textbf{Question.} Each evaluated customer of the public FAR-Trans release~\cite{farTrans,farTransData} gets a Top-10 asset list, scored by binary nDCG@10 against the assets bought in the next 30 days, with past-only inputs~\cite{starter}. The provided \texttt{latest\_buy} list has two parts: a \emph{prefix} with the customer's own eligible past assets, most recently bought first, and a \emph{fill} of never-bought assets in popularity order. My teammate A.~Kosychev studies the fill; I study the prefix. The Session~1 question (Ranker track) was:
\begin{quote}
Does \textbf{zero-shot reranking of a customer's own previously bought assets by an open-weight LLM}, given dated Buy and Sell events and as-of asset metadata, improve \textbf{nDCG@10 on purchases in the next 30 days} over \textbf{a logistic-regression ranker that receives the same information as features}, for \textbf{warm FAR-Trans customers}?
\end{quote}

\textbf{What changed since Session~1.} The LLM arm was planned then and is now implemented and run. After reading the first outputs I added one control that was not planned: the order in which the prompt lists the assets, with no model. Section~\ref{sec:exp} shows why it is needed. To keep the fill of \texttt{latest\_buy} inside the ranker's candidate pool, the starter runs with all eligible assets as candidates (Section~\ref{sec:fair}).

\textbf{Place in the pipeline.} Past events before the cutoff $\rightarrow$ prefix $P$ (own past assets) $\rightarrow$ reorder $P$ (recency, logistic regression, prompt order or LLM) $\rightarrow$ append the same popularity fill $\rightarrow$ Top-10 $\rightarrow$ the starter's nDCG@10. Only the order of $P$ differs between arms, so any difference comes from the reranker.

\section{Related Work}

The nearest prior work is LLMRank~\cite{llmrank}: a zero-shot LLM ranks candidates given the user's interaction sequence. It reports that LLMs perceive the order of the history poorly and show position bias towards the candidates listed first. The course's adaptation~\cite{baselines} reorders 50 CBF candidates, mostly never-bought assets, and is compared with the CBF order. My arm reorders only assets the customer already bought and is compared with a learned ranker with the same facts. Work on repeat consumption shows that repeat items carry much of the accuracy of session and next-basket recommenders~\cite{repeatnet,nbrcheck}; RepeatNet answers the repeat question with a trained decoder, not a frozen LLM. FLARKO~\cite{flarko} fine-tunes an LLM on FAR-Trans with knowledge graphs, but does not isolate repeat ordering or report nDCG@10 on next-month purchases. I did not find a study of zero-shot LLM reordering of a customer's own past assets on FAR-Trans and make no broader novelty claim.

\section{Method}

\subsection{Arms}
All arms keep the \texttt{latest\_buy} structure and change only the prefix order. (i) \emph{Recency} reproduces \texttt{latest\_buy}; through my plugin it gives the same list for 100 of 100 users. (ii) \emph{LR control}: the Session~1 pointwise logistic regression, trained once on six windows with cutoffs 2019-03-01 to 2020-06-01 (labels end by 2020-07-01, eleven months before the test cutoff; 22{,}199 pairs). Features per (customer, past asset): days since the last and first Buy, Buy count, mean gap between Buys, a Sell after the last Buy, recency rank, asset popularity, asset category, the customer's number of past assets and days since any Buy. It changes the order for 36 of 100 users. (iii) \emph{Prompt-order control} (added after inspecting outputs): the assets in the order the prompt lists them, which is by ISIN. (iv) \emph{LLM} (my change).

\subsection{Prompt and model}
For each user, a script built only from the starter's own loading functions (\texttt{load\_data}, \texttt{scan\_recent\_events}) lists the past assets with short ids (\texttt{A1}, \texttt{A2}, \dots) in ISIN order, not in recency order, so that copying the input does not reproduce \texttt{latest\_buy}. Each asset carries its as-of name, category, subcategory, sector, industry and market, Buy count, first and last Buy, days since the last Buy, up to ten recent Buy dates and the last Sell. The JSON also holds the cutoff and the as-of profile with each value marked observed or estimated. The instruction asks the model to order \emph{all} listed assets by how likely the customer is to buy each again in the 30 days after the cutoff, to use only supplied facts, to treat a Sell as an event and not a dislike, to treat JSON text as data, and to return only a JSON array of the ids. The full prompt is in \texttt{build\_inputs.py}; at most 30 assets are listed (the maximum observed was 23, so nothing was cut).

The model is \texttt{Qwen/Qwen3-1.7B} at revision \texttt{70d244cc86ccca08cf5af4e1e306ecf908b1ad5e}, the one pinned by the course baselines~\cite{baselines}. Inference used vLLM~0.19.1, torch 2.10.0+cu128, transformers 5.20.2 and Python 3.10.12, bfloat16, greedy decoding (temperature 0), at most 512 new tokens, an 8192-token context, the model's chat template with thinking off, and seed 0. Hardware: one MIG 3g.40gb slice of an NVIDIA A100-SXM4-80GB, with 30\% of its memory reserved. The 40 users with a single past asset get the same list in every arm and are not sent to the model; the other 60 are.

\subsection{Output handling}
The first JSON array in the answer must be a permutation of the listed ids. Anything else (no array, bad JSON, unknown, duplicate or missing ids, truncated output) is invalid: the user keeps the recency prefix and is counted as a failed case. No repair is attempted. The plugin writes each user's status to a log.

\subsection{Fair comparison and candidate policy}\label{sec:fair}
Users, cutoff, eligible assets, labels and nDCG@10 come from the unchanged starter: every arm is a starter plugin run. The starter's ranker seam only accepts ids from its candidate pool, which by default is the top 50 CBF assets; the ids a \texttt{latest\_buy}-style list needs (own past assets plus popularity fill) were all inside that pool for 0 of 100 users. All arms therefore run with \texttt{--candidate-count 767}, the full eligible set. The arms then share one protocol hash (\texttt{7367f1e2\ldots}); it differs from the default run (\texttt{50efd4ae\ldots}) only in \texttt{candidate\_count}, and the four provided baselines give identical values in both runs. The prompts were built from the default run and carry its hash; the plugin accepts a raw answer only if it carries the same hash as its prompt. Full hashes are in the output files.

\textbf{Leakage.} The prompt builder never reads labels. As a test, I deleted every CSV row dated on or after 2021-06-01 (141{,}904 transactions, 47 asset rows, 7{,}321 profile rows, 211{,}952 prices) and rebuilt the inputs: prompts and prefixes were byte-identical to the originals.

\section{Experiments}\label{sec:exp}

\begin{table*}[t]
\caption{Results on the starter's sample: 100 users, cutoff 2021-06-01, 30 days, repeats allowed. All arms in one protocol (\texttt{candidate\_count} 767)}
\label{tab:main}
\centering
\footnotesize
\setlength{\tabcolsep}{4pt}
\begin{tabular}{@{}lrrrr>{\raggedright\arraybackslash}p{4.9cm}>{\raggedright\arraybackslash}p{1.9cm}@{}}
\toprule
Method & Users & Users with & nDCG@10 & nDCG@10 & Runtime & Failed\\
 & & future purchases & all users & nonempty truth & & cases\\
\midrule
Popularity & 100 & 11 & 0.0286 & 0.2599 & 6.0 s, one starter run for all four (CPU) & 0\\
CBF & 100 & 11 & 0.0201 & 0.1827 & same run & 0\\
Latest buy & 100 & 11 & 0.0277 & 0.2520 & same run & 0\\
Latest buy or sell & 100 & 11 & 0.0271 & 0.2464 & same run & 0\\
\midrule
LR control (same facts) & 100 & 11 & \textbf{0.0343} & \textbf{0.3116} & fit 12.6 s + scoring 6.1 s, each with data loading (CPU) & 0\\
Prompt-order control (no model) & 100 & 11 & 0.0341 & 0.3103 & not timed separately & 0\\
\textbf{LLM, Qwen3-1.7B (mine)} & 100 & 11 & 0.0316 & 0.2869 & load 57.5 s + generation 6.4 s (A100 slice) & 3 invalid of 60 calls\\
\bottomrule
\end{tabular}
\end{table*}

\textbf{Main result.} Table~\ref{tab:main} gives all arms. Input building takes 5.9~s and each plugin run of the starter 6--11~s on a laptop CPU, the four baselines included. The 60 prompts held 71{,}265 tokens; the answers held 1{,}657 tokens (mean 27.6, maximum 110), and all 60 ended with a stop token. Against \texttt{latest\_buy}, the LLM is better for 3 users, worse for 3 and equal for 94. The LR control is better for 5 and worse for 1. The prompt-order control is better for 3 and worse for 2.

\textbf{What the LLM returns.} Of the 57 valid answers, 46 are exactly the order in which the prompt listed the assets. Only 17 equal the recency order, and these include short lists where the two orders coincide. The model mostly echoes its input and does not use the dates. This matches the position bias reported for LLMRank~\cite{llmrank}. The prompt-order control, which does the same without a model, scores 0.0341, above the LLM; against it the LLM is better for 1 user and worse for 2.

\begin{table}[t]
\caption{The 11 users with purchases (ids truncated). nDCG@10 per arm}
\label{tab:users}
\centering
\footnotesize
\setlength{\tabcolsep}{4pt}
\begin{tabular}{@{}lrrlrrrr@{}}
\toprule
User & Past & Bought & LLM & Rec. & LR & Prompt & LLM\\
 & assets & & status & & & order & \\
\midrule
29C0 & 20 & 1 & invalid & 0 & 0 & 0 & 0\\
6FA5 & 18 & 6 & valid & 0.4725 & 0.6448 & 0.2874 & 0.2874\\
8A22 & 17 & 1 & invalid & 0 & 0 & 0 & 0\\
94A1 & 14 & 1 & valid & 0.3155 & 0.4307 & 0.3010 & 0\\
5180 & 13 & 2 & valid & 0 & 0.2372 & 0.1846 & 0.1846\\
6BAA & 6 & 1 & valid & 0.3869 & 0.5000 & 0.3869 & 0.5000\\
D63D & 6 & 1 & valid & 0.3333 & 0.3333 & 0.3333 & 0.3333\\
870E & 5 & 2 & valid & 0.8772 & 0.8503 & 0.9197 & 0.8503\\
5087 & 5 & 1 & valid & 0.3869 & 0.4307 & 1.0000 & 1.0000\\
0E72 & 4 & 1 & valid & 0 & 0 & 0 & 0\\
6B15 & 2 & 1 & valid & 0 & 0 & 0 & 0\\
\bottomrule
\end{tabular}
\end{table}

\textbf{Who moves.} Table~\ref{tab:users} lists the users with purchases. Customer 5087 is the data example of Session~1: the purchased asset happens to have the smallest ISIN of the customer's five, so any arm that keeps the ISIN order puts it first (nDCG 1.0 instead of 0.3869). This one customer adds 0.0061 to the all-user mean of the LLM; without him the LLM would be 0.0023 \emph{below} \texttt{latest\_buy}. Where the LLM departs from its input order, it helps once (6BAA) and hurts twice (870E, and 94A1, whose purchased asset drops out of the top ten). The LR control is better than recency for five of these users and worse for one (870E).

\textbf{Invalid answers.} All three invalid answers dropped ids: for a user with six listed assets the model returned \texttt{["A1", "A3", "A4", "A5", "A6"]}, and for users with 20 and 17 assets it returned 16 and 14 ids. Two of the three are labelled users; every arm, the fallback included, scores 0 for them.

\textbf{Repeatability.} A second greedy pass over the same 60 prompts in the same process changed 5 raw answers. The likely cause is that vLLM batches requests and reuses cached prefixes, so the floating-point path differs between passes; I did not test this further. Only the first pass is evaluated; its raw text is saved and the plugin reproduces the reported lists from it.

\section{Discussion}

\textbf{What improved and what failed.} The pipeline works end to end: a pinned open-weight model, a past-only input builder, an evaluation through the unchanged starter and a logged fallback. The research answer on this sample is negative. The LLM is above \texttt{latest\_buy}, but below the learned control with the same facts and below a model-free control that copies its input order. The gain over \texttt{latest\_buy} comes from one customer whose asset sorts first by ISIN. With 11 labelled users a single user moves the all-user mean by up to 0.006, so no ranking of methods is claimed. For the product decision framed in Session~1, nothing here argues for per-customer LLM inference over the logistic-regression order.

\textbf{A failure during the work.} After the LLM run, the analysis script crashed:
\begin{quote}\footnotesize\ttfamily
lr, log = jl(INP / 'lr\_orders.jsonl'),\\
\hspace*{1em}jl(INP / 'plugin\_log\_llm.jsonl')\\
FileNotFoundError: [Errno 2] No such file or\\
\hspace*{1em}directory: '...\textbackslash plugin\_log\_llm.jsonl'
\end{quote}
Root cause: the analysis imported the plugin to reuse its parser, and the plugin deleted its own log at import time, so the analysis erased the log it was about to read. Fix: the parser moved into a module without side effects (\texttt{llm\_parse.py}), and the plugin now resets its log on its first call, not on import. Verification: all arms were rerun, each log has 100 rows, and the metrics did not change (LLM 0.0316, LR 0.0343). A second, smaller failure: PowerShell~5.1 read the runner script without a byte-order mark as ANSI and garbled the Cyrillic folder names; the script now builds its paths from its own location.

\textbf{Next controlled experiment.} \emph{Changes:} only the order in which the prompt lists the past assets: ISIN order is replaced by three seeded random permutations, with ids reassigned. \emph{Fixed:} model, revision, decoding, instruction text, users, fill, metric and the strict output policy. \emph{Measured:} the share of answers equal to the listing order, the rank correlation with recency, and nDCG@10. If the answers follow the listing order whatever it is, the model does not read dates, and the prompt-order control is the right comparison. The test moves to the nine quarterly development cutoffs of Session~1 (12{,}163 labelled user-windows), as Homework~3 allows, because 11 users cannot separate the arms.

\textbf{Threats to validity.} One model, one prompt and one sample with 11 labelled users. The prompt-order control was added after the outputs were seen. The ranker runs with all eligible assets as candidates, not the default 50. Greedy decoding in vLLM was not exactly repeatable. No suitability, return or profit claim is made.

\textbf{Team discussion and individual contribution.} Kosychev pushed his Session~2 solution to the team repository on 25 November (folder \texttt{hw-module-2-2}): frozen Qwen3-Embedding-0.6B vectors of asset text order the fill behind the unchanged \texttt{latest\_buy} prefix. On the same 100 users it is 0.0017 above \texttt{latest\_buy}; on the nine development cutoffs it is the weakest fill for users who buy an asset new to them. Our two changes touch disjoint parts of the list, so the baseline / prefix-only / fill-only / both comparison proposed in Session~1 stays possible. His draft records the process the team agreed for this homework: each member first builds an individual solution, and the solutions are merged afterwards. At the time of writing the merge has not taken place, so no joint decision on the combined pipeline exists yet. Open questions: which prefix and fill enter the team's benchmark submission, whether the development protocol replaces the 100-user sample for decisions, and how the work of Patrushev and Vaal joins the pipeline. \emph{My contribution:} everything in \texttt{hw-module-2-2/starchenko}: the input builder and prompt, the inference on the A100, the plugin and output policy, the LR control (my Session~1 code, unchanged), the prompt-order control, the leakage test and the analysis. \emph{Shared:} the course starter and Kosychev's nine-cutoff protocol.

\section{AI Usage Disclosure}

\textbf{Tools.} Claude Code (Claude Opus~5.5, \texttt{claude-opus-5-5}, VS Code on Windows) wrote all code, ran the generation on the A100 over ssh (my choice of hardware instead of Kaggle), ran the starter and the analysis, committed the code and drafted this report.

\textbf{What the AI checked.} The plugin's recency arm equals \texttt{latest\_buy} for 100 of 100 users; the prompt builder's prefixes equal the starter's lists; a rebuild of the inputs is byte-identical; the truncation test above; the two \texttt{protocol.json} files differ only in \texttt{candidate\_count}; the baselines equal the course's published HW2 results and Session~1; the LR weights equal those of Session~1.

\textbf{What I personally verified.} (1)~For customer \texttt{5087D23F}, I took the truth from the raw \texttt{transactions.csv}: one purchase in June 2021, \texttt{GRS003003035} on 2021-06-08. In \texttt{predictions.csv} this asset is fifth in the \texttt{latest\_buy} list, so $1/\log_2 6 = 0.3869$, and first in the LLM list, so $1/\log_2 2 = 1.0$; \texttt{per\_user.csv} holds the same values. (2)~For the same customer I mapped the raw answer \texttt{["A1","A2","A3","A4","A5"]} through \texttt{alias\_to\_isin}: the five ISINs are exactly the first five positions of the LLM list, and positions 6--10 equal the tail of \texttt{latest\_buy}. The answer repeats the prompt's order. In all 60 prompts the ids follow alphabetical ISIN order, and of the 60 answers I counted 46 that repeat this order, 11 that change it and 3 invalid ones. The 1.0 for this customer is therefore an accident of the purchased asset sorting first by ISIN. (3)~For customer \texttt{F3B6B4C7}, the prompt lists six assets and the model returned five ids, without \texttt{A2}; the plugin log shows \texttt{fallback\_missing\_ids}, and the LLM and \texttt{latest\_buy} lists in \texttt{predictions.csv} agree in all ten positions. Over the whole log there are 57 valid answers, 40 users without a call and 3 fallbacks. (4)~In the prompt of customer 5087, 2021-06-01 appears once, as the cutoff field; all other dates are earlier (the earliest is 2020-01-15), and the June purchase does not appear. None of the 60 prompts contains a date after 2021-06-01. (5)~The \texttt{protocol.json} files of the default run and the LLM run differ in two fields, \texttt{candidate\_count} (50 against 767) and \texttt{protocol\_hash} (\texttt{50efd4ae\ldots} against \texttt{7367f1e2\ldots}); the prompts carry the hash of the default run. The four baselines of the default run match the LLM run's console output after rounding to four decimals, and that output gives 0.0316 / 0.2869 for the LLM.

\textbf{What I trusted without verification.} I did not rerun the GPU generation, the truncation test or the arm runs myself; check~(4) covers the dates in the prompts, not every input path. I also relied on the vLLM chat template and tokenizer handling, the course starter's eligibility rules, and the assumption that the five answers that changed in the repeat pass would not change the conclusion.

\textbf{Session export.} The unmodified Claude Code JSONL of this session, submitted as \texttt{starchenko\_session2\_\allowbreak session.json}.

\begin{thebibliography}{9}
\bibitem{farTrans} J.~Sanz-Cruzado, N.~Droukas, and R.~McCreadie, ``FAR-Trans: An investment dataset for financial asset recommendation,'' in \emph{Proc. IJCAI-24 Workshop on Recommender Systems in Finance (Fin-RecSys)}, 2024. [Online]. Available: \url{https://arxiv.org/abs/2407.08692}
\bibitem{farTransData} University of Glasgow, ``FAR-Trans dataset,'' research data repository. [Online]. Available: \url{https://researchdata.gla.ac.uk/1658/}
\bibitem{repeatnet} P.~Ren, Z.~Chen, J.~Li, Z.~Ren, J.~Ma, and M.~de~Rijke, ``RepeatNet: A repeat aware neural recommendation machine for session-based recommendation,'' in \emph{Proc. AAAI Conf. Artificial Intelligence}, 2019. [Online]. Available: \url{https://arxiv.org/abs/1812.02646}
\bibitem{nbrcheck} M.~Li, S.~Jullien, M.~Ariannezhad, and M.~de~Rijke, ``A next basket recommendation reality check,'' \emph{ACM Trans. Information Systems}, vol.~41, no.~4, 2023. [Online]. Available: \url{https://arxiv.org/abs/2109.14233}
\bibitem{llmrank} Y.~Hou, J.~Zhang, Z.~Lin, H.~Lu, R.~Xie, J.~McAuley, and W.~X.~Zhao, ``Large language models are zero-shot rankers for recommender systems,'' in \emph{Proc. European Conf. Information Retrieval (ECIR)}, 2024. [Online]. Available: \url{https://arxiv.org/abs/2305.08845}
\bibitem{flarko} F.~Spadea and O.~Seneviratne, ``Aligning language models with investor and market behavior for financial recommendations,'' arXiv:2510.15993, 2025. [Online]. Available: \url{https://arxiv.org/abs/2510.15993}
\bibitem{starter} S.~Jin, ``FAR-Trans student starter,'' LLM4Rec course repository, 2026. [Online]. Available: \url{https://github.com/dryjins/RecSys-LLMs/tree/main/project/far-trans}
\bibitem{baselines} S.~Jin, ``FAR-Trans: Encoder and Ranker baselines,'' LLM4Rec course repository, 2026. [Online]. Available: \url{https://github.com/dryjins/RecSys-LLMs/tree/main/project/far-trans-llm}
\end{thebibliography}

\end{document}
